%!TEX root =  tfg.tex
\chapter{Fase 2: Diseño e Implementación}

\begin{abstract}
Este capítulo recoge el trabajo realizado durante la Fase 2 del proyecto, organizada por features
del FDD. Para cada feature se explica su objetivo general y, cuando ha sido necesario, su desglose
en varias issues de backlog distribuidas en distintas iteraciones.
\end{abstract}

\section{Objetivos de la fase}
Durante esta fase se definió e implementó la base funcional del videojuego. El objetivo principal
fue construir los sistemas esenciales sobre los que se apoya el resto del proyecto, garantizando
que la experiencia de juego pudiera desarrollarse de forma estable y coherente desde el inicio.

En esta etapa se abordaron tanto la generación del mundo como la estructura básica de juego,
la interacción del jugador con el entorno y la interfaz mínima necesaria para que la partida
resultara jugable. El seguimiento se ha realizado a nivel de feature general, vinculando su
ejecución a issues concretas cuando la complejidad obligó a dividir el trabajo en varias iteraciones.

\section{Desarrollo por features}
El trabajo realizado durante esta fase siguió el flujo de trabajo definido para el proyecto y
se ejecutó de forma incremental. La planificación se controló por features del FDD y, dentro de
cada feature, por issues de backlog que permitieran cerrar entregables verificables por iteración.

\subsection{Feature: Generación Procedural}
Esta feature define el sistema que construye el nivel de forma dinámica. Incluye la estructura del
mapa, la conexión entre zonas y las reglas de distribución del contenido inicial.

\begin{enumerate}
\item \textbf{Issue: Generación Procedural (I) -- Iteración 1.} Se implementó una primera versión funcional para disponer de una base jugable donde validar movimiento, aparición de entidades y navegación inicial.
\item \textbf{Issue: Generación Procedural (II) -- Iteración 2.} Se completó el refinamiento del sistema, ajustando distribución y estabilidad de la estructura del nivel.
\end{enumerate}

Esta feature se dividió en dos issues porque su complejidad real fue mayor de la estimada al inicio.
La primera entrega permitió validar el enfoque técnico, pero no era suficiente para dar por cerrada
la feature completa sin una segunda iteración de ajustes.

\subsection{Feature: Lógica de Círculos}
Esta feature gestiona la progresión del jugador por niveles temáticos y su coherencia con la estructura
de partida.

\begin{enumerate}
\item \textbf{Issue: Lógica de Círculos (I) -- Iteración 2.} Se construyó la versión inicial del flujo de progresión.
\item \textbf{Issue: Lógica de Círculos (II) -- Iteración 3.} Se revisó y consolidó la lógica tras validar su integración con la generación del nivel.
\end{enumerate}

Se separó en dos issues porque la progresión dependía de componentes previos todavía inestables en la
iteración 2. La segunda issue permitió cerrar inconsistencias detectadas durante la integración.

\subsection{Feature: Minimap}
Esta feature proporciona apoyo de navegación al jugador mediante una representación simplificada del nivel.

\begin{enumerate}
\item \textbf{Issue: Minimap (I) -- Iteración 3.} Se implementó la base visual y funcional para consulta de orientación.
\item \textbf{Issue: Minimap (II) -- Iteración 6.} Se completó el ajuste final cuando la estructura completa del juego ya estaba estabilizada.
\end{enumerate}

La división en dos issues se debió a que el minimapa depende directamente del estado final de salas,
rutas y sistemas de progresión. Cerrar la feature en una sola iteración habría generado retrabajo.

\subsection{Feature: Menús}
Esta feature engloba la navegación de pantallas y los flujos de entrada, pausa y salida del juego.

\begin{enumerate}
\item \textbf{Issue: Menús (I) -- Iteración 4.} Se implementaron los flujos principales de navegación.
\item \textbf{Issue: Menús (II) -- Iteración 4.} Se refinó el comportamiento de transición y consistencia del flujo de pantallas.
\end{enumerate}

Aunque ambas issues se resolvieron en la misma iteración, se separaron para mantener trazabilidad entre
la construcción base del sistema y su posterior consolidación de uso.

\subsection{Features no divididas en iteraciones}
El resto de trabajo de la fase se gestionó con una sola issue por feature o por bloque funcional:

\begin{itemize}
\item \textit{Añadir Jugador} (Iteración 1).
\item \textit{Añadir Enemigo} (Iteración 1).
\item \textit{Cámara del Jugador} (Iteración 2).
\item \textit{Dualidad de Personajes} (Iteración 3).
\item \textit{Head-Up Display (HUD)} (Iteración 4).
\item \textit{Actualizar Sprites (I)} (Iteración 5).
\item \textit{Sistema de Combate (I)} (Iteración 5).
\item \textit{Sistema de Enemigos (I)} (Iteración 5).
\item \textit{Gestión de Inventario} (Iteración 6).
\item \textit{Objetos} (Iteración 6).
\end{itemize}

\section{Validación y resultados}
La validación de esta fase se centró en comprobar que todos los sistemas núcleo quedaban integrados
entre sí y que el jugador podía avanzar dentro de una versión funcional del proyecto. Se realizaron
pruebas sobre la generación de niveles, el comportamiento del personaje, la navegación entre salas
y la correcta visualización de la interfaz.

Como resultado, quedó consolidada una base jugable sobre la que fue posible seguir desarrollando
contenido y ampliar el proyecto en fases posteriores.

\section{Incidencias y ajustes}
Durante la fase se realizaron ajustes sobre la \gls{generacionprocedural}, la cámara y la estructura de
las \glspl{sala} para asegurar que todos los sistemas funcionaran de forma coordinada. También fue necesario
adaptar distintos elementos de la interfaz para que la información del juego se mostrara de forma clara
y consistente.

En concreto, se revisaron los sistemas de \gls{hud}, \gls{minimapa}, \gls{inventario}, \gls{objeto}, \gls{menu} y \gls{combate} para que el cierre de cada iteración no dejara tareas abiertas que pertenecieran a la siguiente. Este criterio evitó que una issue quedara artificialmente cerrada antes de tiempo.

Estos ajustes permitieron estabilizar el núcleo del juego y dejaron preparada la arquitectura base
para la incorporación de nuevas funcionalidades.

\section{Resumen de horas de la fase}

\begin{table*}[htb]
	\centering
	\begin{coolTable}{ll}{2}
{Horas de la fase 2}
	\textbf{Horas previstas}&Pendiente de ajustar\\
	\textbf{Horas reales}&158\\
	\textbf{Observaciones}&\\
	\end{coolTable}
	\caption{Resumen de horas de la Fase 2}
\end{table*}
